% ===========================================================================
%  UIU GPU server - student guide
%  All addresses, usernames and passwords are PLACEHOLDERS.
%
%  Colour code:
%    amber  = type this on your Windows laptop
%    green  = type this on the server
% ===========================================================================
\documentclass[12pt,a4paper]{article}

\usepackage[T1]{fontenc}
\usepackage[utf8]{inputenc}
\usepackage[default]{lato}
\usepackage[varqu,varl]{inconsolata}
\usepackage[margin=1.5cm,top=1.2cm,bottom=1.2cm]{geometry}
\usepackage{xcolor}
\usepackage{colortbl}
\usepackage{titlesec}
\usepackage{enumitem}
\usepackage{tabularx}
\usepackage{booktabs}
\usepackage[most]{tcolorbox}
\tcbuselibrary{listings,skins}
\usepackage{needspace}
\usepackage[hidelinks]{hyperref}

\definecolor{ink}{HTML}{16211C}
\definecolor{forest}{HTML}{1F5134}
\definecolor{amber}{HTML}{8A4B00}
\definecolor{rose}{HTML}{9F1239}
\definecolor{teal}{HTML}{0F766E}
\definecolor{linkc}{HTML}{1D4ED8}

\hypersetup{colorlinks=true, urlcolor=linkc, linkcolor=linkc}
\linespread{1.06}

\titleformat{\section}
  {\large\bfseries\color{forest}}{\textcolor{teal}{\thesection}}{0.5em}{}
  [{\vspace{-7pt}\noindent{\color{forest!30}\rule{\linewidth}{0.8pt}}}]
\titlespacing*{\section}{0pt}{10pt}{6pt}

\setlength{\parindent}{0pt}
\setlength{\parskip}{3pt}

\tcbset{
  cmdstyle/.style={
    enhanced, boxrule=0pt, arc=3pt,
    leftrule=3.5pt, left=10pt, right=10pt, top=4pt, bottom=4pt,
    before skip=4pt, after skip=4pt,
  }
}

\newtcblisting{winblock}{
  cmdstyle, colback=amber!12, colframe=amber!85,
  listing only,
  listing options={
    basicstyle=\ttfamily\small\color{ink},
    commentstyle=\color{amber!90}\itshape,
    columns=fullflexible, keepspaces=true, breaklines=true,
    aboveskip=0pt, belowskip=0pt, xleftmargin=0pt, xrightmargin=0pt,
  }
}

\newtcblisting{srvblock}{
  cmdstyle, colback=forest!12, colframe=forest!85,
  listing only,
  listing options={
    basicstyle=\ttfamily\small\color{ink},
    commentstyle=\color{forest!90}\itshape,
    columns=fullflexible, keepspaces=true, breaklines=true,
    aboveskip=0pt, belowskip=0pt, xleftmargin=0pt, xrightmargin=0pt,
  }
}

\newcommand{\cmd}[1]{{\ttfamily\small #1}}

\newtcolorbox{titleband}{colback=forest, colframe=forest, boxrule=0pt, arc=4pt,
  left=12pt, right=12pt, top=10pt, bottom=10pt, halign=center}
\newtcolorbox{quick}{colback=teal!8, colframe=teal!80, boxrule=0.9pt, arc=4pt,
  left=8pt, right=8pt, top=6pt, bottom=6pt, fontupper=\small,
  title={\small\bfseries Quick start}, coltitle=white, colbacktitle=teal!85}
\newtcolorbox{tipbox}{colback=teal!6, colframe=teal!55, boxrule=0.7pt, arc=3pt,
  left=8pt, right=8pt, top=5pt, bottom=5pt, fontupper=\small}
\newtcolorbox{warnbox}{colback=rose!6, colframe=rose!70, boxrule=0.7pt, arc=3pt,
  left=8pt, right=8pt, top=5pt, bottom=5pt, fontupper=\small}
\newtcolorbox{legend}{colback=white, colframe=forest!30, boxrule=0.7pt, arc=3pt,
  left=8pt, right=8pt, top=5pt, bottom=5pt, fontupper=\small}

\begin{document}

\begin{titleband}
{\LARGE\bfseries\color{white} How to use the shared GPU server}\\[3pt]
{\small\color{white} Ubuntu 24.04 Server \textperiodcentered\ SSH \textperiodcentered\ 200 GB \textperiodcentered\ monitored by UIU CITS}
\end{titleband}

\begin{warnbox}
\textbf{Every address, username and password below is a placeholder.} Replace them with the values
issued to you, and never share real credentials.
\end{warnbox}

\begin{quick}
\begin{enumerate}[leftmargin=1.2em,itemsep=1pt,topsep=0pt]
  \item Save the short name once (section 3), then connect with \cmd{ssh uiu-gpu}.
  \item Install \textbf{VS Code} with the \textbf{Remote - SSH} extension (section 4).
  \item Queue every job with \cmd{ts} so the shared GPU is not overloaded (section 6).
\end{enumerate}
\end{quick}

\begin{legend}
\textbf{Two colours, two places.}\quad
\colorbox{amber!20}{\textbf{\textcolor{amber}{amber}}}~type it on your Windows laptop.\quad
\colorbox{forest!20}{\textbf{\textcolor{forest}{green}}}~type it on the server.
\end{legend}

\section{What you are given}

\begin{tabularx}{\linewidth}{@{}>{\bfseries}lX@{}}
\toprule
\rowcolor{forest!10} Item & Detail \\
\midrule
Operating system & Ubuntu 24.04 Server, no desktop \\
Access & SSH only; there is no remote desktop \\
File transfer & over the same SSH connection, using \cmd{scp} \\
Storage & 200 GB total; check with \cmd{df -h} on the server \\
GPUs & shared, so keep one job running at a time \\
Support & monitored at all times by UIU CITS \\
\bottomrule
\end{tabularx}

\section{Connecting from a terminal}

Windows 10 and 11 already include the SSH client. Check it once:

\begin{winblock}
ssh -V
\end{winblock}

If it prints a version you are ready. If not, add \emph{OpenSSH Client} from
\emph{Settings $\rightarrow$ Apps $\rightarrow$ Optional features}, then open a new terminal.

\begin{winblock}
ssh username@10.100.110.110          # on the campus network
ssh -p 110 username@100.100.52.109   # from outside, through the port
\end{winblock}

Accept the host key with \cmd{yes} on the first connection. Nothing appears while you type the
password, which is normal. When the prompt changes, you are on the server.

\section{Save a short name: \texttt{ssh uiu-gpu}}

Typing that address every time is slow and easy to mistype. Windows keeps a small file that
\cmd{ssh}, \cmd{scp} and VS Code all read, so you can save the address once and then use a short
name forever. This works in PowerShell and in Command Prompt.

\textbf{Step 1.} Create the folder, if it does not exist yet. Run \textbf{one} of these:

\begin{winblock}
mkdir "%USERPROFILE%\.ssh"                         # Command Prompt
New-Item -ItemType Directory -Force "$HOME\.ssh"   # PowerShell
\end{winblock}

\textbf{Step 2.} Open the file in Notepad:

\begin{winblock}
notepad "%USERPROFILE%\.ssh\config"
\end{winblock}

If Notepad asks whether to create the file, say yes.

\textbf{Step 3.} Paste this, put your own account name in place of \cmd{username}, and save:

\begin{winblock}
Host uiu-gpu
    HostName 10.100.110.110
    User username
    Port 22
\end{winblock}

\textbf{Step 4.} Close the file and test it:

\begin{winblock}
ssh uiu-gpu
\end{winblock}

That is all. From now on \cmd{ssh uiu-gpu}, \cmd{scp} and VS Code all know the address.
From outside the campus network, add a second block and use \cmd{uiu-gpu-outside}:

\begin{winblock}
Host uiu-gpu-outside
    HostName 100.100.52.109
    User username
    Port 110
\end{winblock}

\begin{warnbox}
Two things go wrong for beginners. The file must be called exactly \cmd{config}, with no
\cmd{.txt} after it; in Notepad choose \emph{Save as} $\rightarrow$ \emph{All files}. And keep the
four lines on separate lines with the indentation shown --- do not join them into one line.
\end{warnbox}

\section{Connecting with VS Code (recommended)}

VS Code gives you a file tree, a terminal and an editor on the server. It is far easier than typing
every command, and you can drag files in and out.

\begin{enumerate}[leftmargin=1.3em,itemsep=1.5pt,topsep=2pt]
  \item Install \textbf{Visual Studio Code} from \url{https://code.visualstudio.com}.
  \item Open Extensions (\cmd{Ctrl+Shift+X}) and install \textbf{Remote - SSH} by Microsoft.
  \item Press \cmd{F1}, type \cmd{Remote-SSH: Connect to Host}, then choose \cmd{uiu-gpu}
        (it appears automatically once section 3 is done).
  \item Choose \emph{Linux}, then \emph{Open Folder}, then \cmd{~/work}.
  \item Use \cmd{Ctrl+\`{}} for a terminal. It already runs on the server.
\end{enumerate}

\begin{tipbox}
Upload and download by dragging files into the VS Code sidebar, or right-click a file and choose
\emph{Download}. It is the same transfer as \cmd{scp}, just with a mouse. Full guide:
\url{https://code.visualstudio.com/docs/remote/ssh}
\end{tipbox}

\section{Moving files with \texttt{scp}}

Run these on your laptop, in a \textbf{new} PowerShell window, not inside the SSH session. The
short name from section 3 works here too.

\begin{winblock}
scp .\train.py username@10.100.110.110:~/work/   # full address
scp .\train.py uiu-gpu:~/work/                   # short name, same result
scp -r .\project uiu-gpu:~/work/                 # upload a folder
scp uiu-gpu:~/work/out.csv .                     # download a file
scp -r uiu-gpu:~/work/results .                  # download a folder
\end{winblock}

Without the short name, port forwarding uses a capital \cmd{-P}:

\begin{winblock}
scp -P 110 .\train.py username@100.100.52.109:~/work/
\end{winblock}

Capital \cmd{-P} belongs to \cmd{scp}; lower-case \cmd{-p} belongs to \cmd{ssh}. Full manual:
\url{https://man.openbsd.org/scp}

\section{Running jobs with Task Spooler}

The GPU is shared, so never leave training in a plain terminal. Queue every job instead. Task
Spooler runs your jobs one after another and keeps each output in its own file, so you can close
the laptop and read the result later.

\begin{srvblock}
ts python train.py     # add a job and return immediately
ts                     # list the queue: ID, state, time, command
ts -t 3                # follow job 3 live (tail)
ts -c 3                # print all output of job 3 (cat)
ts -k 3                # kill job 3
ts -C                  # clear finished jobs from the list
\end{srvblock}

A typical session:

\begin{srvblock}
ts -S 1                                        # one job at a time
ts bash -c "source ~/venv/bin/activate && python train.py --epochs 5"
ts                                             # note the ID, say 7
ts -t 7                                        # watch it, then Ctrl+C to stop watching
\end{srvblock}

More detail, with every flag and example:
\url{https://justanhduc.github.io/2021/02/03/Task-Spooler.html} and
\url{https://github.com/justanhduc/task-spooler}.

\section{Watching the machine}

\begin{srvblock}
nvidia-smi              # one snapshot: memory, load, processes
watch -n 2 nvidia-smi   # refresh every 2 seconds
df -h                   # disk space
du -sh ~/work           # size of your workspace
\end{srvblock}

Check that the GPU is idle before you start. If \cmd{nvidia-smi} shows a process that is not yours,
wait or ask. Two jobs on one GPU are slower than one, and timings become meaningless.

\section{Python environments}

Work in your own environment; never install into the system Python.

\begin{srvblock}
python3 -m venv ~/venv
source ~/venv/bin/activate
pip install --upgrade pip && pip install -r requirements.txt
\end{srvblock}

\section{Rules you must follow}

\begin{warnbox}
\begin{itemize}[leftmargin=1.1em,itemsep=1pt,topsep=0pt]
  \item Avoid installing unnecessary software packages.
  \item Back up all necessary data every day.
  \item Do not share your user credentials with anyone.
  \item The server is monitored at all times by UIU CITS.
\end{itemize}
\end{warnbox}

A backup means copying results \textbf{to your own computer}, not trusting the server disk. Download
them at the end of every working day, with VS Code or with \cmd{scp}.

\Needspace{12\baselineskip}
\section{When something goes wrong}

\begin{tabularx}{\linewidth}{@{}>{\ttfamily\small}lX@{}}
\toprule
\rowcolor{rose!10} Message & What to do \\
\midrule
ssh not recognised & the Windows SSH client is missing; install OpenSSH Client (section 2) \\
Permission denied & the account name is wrong, or your key is not on the server; ask the administrator \\
Connection timed out & you are outside the campus network; use \cmd{uiu-gpu-outside} or \cmd{-p 110} \\
\cmd{uiu-gpu} unknown & the \cmd{config} file is missing, named \cmd{config.txt}, or in the wrong folder (section 3) \\
No space left & check \cmd{df -h}, delete old checkpoints, download results and remove them \\
CUDA out of memory & another job holds the GPU, or the batch is too large; check \cmd{nvidia-smi} \\
job looks stuck & \cmd{ts} to see its state, \cmd{ts -t <id>} to see where it stopped \\
\bottomrule
\end{tabularx}

\section{Ubuntu commands you will need}

\begin{srvblock}
ls -lh              # list files with sizes
cd ~/work           # change folder
mkdir results       # make a folder
mv old new          # rename or move
rm old.csv          # delete a file (no undo)
less log.txt        # scroll a file, q to quit
tail -f log.txt     # watch a file grow, Ctrl+C to stop
grep ERROR log.txt  # find lines containing a word
nano notes.txt      # editor: Ctrl+O save, Ctrl+X exit
du -sh *            # size of each item here
\end{srvblock}

\begin{tipbox}
\textbf{Before you log off:} check for jobs you forgot with \cmd{ts}, download anything worth
keeping, and free space with \cmd{df -h} if you are near the 200 GB limit.
\end{tipbox}

\end{document}
